\subsection{A Frobenius-eigenvalue obstruction for the actual quadratic source coefficient}

\begin{proposition}[P-F66: the split coefficient is not a geometric mixed Frobenius eigenvalue]
\label{prop:nonstandardweileigenvalue}
Retain the exact source normalization from P-F39 and P-F52.  Thus
\[
\begin{split}
 A(q)&=5+8q^4-4q^6+18q^8-4q^{10}+8q^{12}+5q^{16},\\
 D(q)&=q^{24}(1+q^8)^2A(q),\\
 F(q)&=q^{-16}(1+q+\cdots+q^9)^2(1+q^2)(1+q^{12}),\\
 P(q)&=q^{20}T(q),
\end{split}
                                                        \tag{P-F66a}
\]
where
\[
\begin{split}
 T(q)={}&22-2(q^2+q^{-2})+6(q^4+q^{-4})
       +2(q^6+q^{-6})\\
       &\hspace{12mm}+3(q^8+q^{-8})+2(q^{10}+q^{-10}).
\end{split}
\]
The two conjugate eigenvalues of the actual prescribed central action
on the descended quadratic source module are
\[
       \alpha_\pm(q)
         =\frac{F(q)\bigl(P(q)\pm\sqrt{D(q)}\bigr)}{2q^{20}}.
                                                        \tag{P-F66b}
\]
They are the negatives of the normalized source coefficient
$a=Fc$ and its conjugate, not newly chosen spectral values.

Fix any $e\ge1$ and $q_0=3^e$.  Then
$\alpha_+(q_0),\alpha_-(q_0)$ are positive, distinct quadratic
conjugates.  For every $b\in\mathbb Q^\times$ and every nonzero
integer $r$, neither
\[
                         b\,\alpha_\pm(q_0)^r
                                                        \tag{P-F66c}
\]
is a Weil number of any weight relative to any prime power $Q$.
Here a \emph{Weil number} means an algebraic number all of whose
complex conjugates have the same absolute value $Q^{w/2}$,
with $w\in\mathbb Z$; integrality is not being imposed.

Consequently none of the numbers in (P-F66c) can occur as a
geometric Frobenius eigenvalue in
\[
           H_c^i(X_{\overline{\mathbb F}_Q},\mathcal L),
                                                        \tag{P-F66d}
\]
where $X/\mathbb F_Q$ is separated of finite type,
$\ell\ne\operatorname{char}\mathbb F_Q$, and $\mathcal L$
is a constructible $\overline{\mathbb Q}_\ell$-sheaf mixed
in Deligne's all-complex-conjugates sense.  This includes
constant coefficients and geometric coefficient constructions
for which that mixedness has been established.  It also
excludes their occurrence in Frobenius-stable subquotients,
finite direct sums, or Tate twists of these groups.
\end{proposition}

\begin{proof}
The exact identifications (P-F66a)--(P-F66b) were established
from GCT VII, Figure~2, printed p.~31, and Figure~7, p.~36,
and the normalization in Section~7.1.7, pp.~27--28:
\url{https://arxiv.org/pdf/0709.0749}.
The coefficient and its minimal-polynomial normalization are
the actual GCT VIII calculation in Sections~5.3--5.4.1,
pp.~31--32 and p.~48:
\url{https://arxiv.org/pdf/0709.0751}.
P-F52 identifies (P-F66b) with the central word $FZ$ in
the actual source representation.  P-F39 proves, for every
positive real $q$,
\[
                 F(q)>0,\qquad D(q)>0,\qquad
                         P(q)>\sqrt{D(q)}.
                                                        \tag{P-F66e}
\]
In particular both numbers in (P-F66b) are positive and
unequal at $q=q_0$.

The specialization is genuinely quadratic.  Since $3\mid q_0$,
\[
                         A(q_0)\equiv5\equiv2\pmod3.
                                                        \tag{P-F66f}
\]
An integer that is a square in $\mathbb Q$ is a square
in $\mathbb Z$: if $(u/v)^2$ is an integer with
$u,v$ relatively prime and $v>0$, then $v^2\mid u^2$
forces $v=1$.  Integer squares modulo $3$ are $0$ and $1$.
Thus $A(q_0)$ is not a square in $\mathbb Q$.
Because
\[
                 D(q_0)=\bigl(q_0^{12}(1+q_0^8)\bigr)^2A(q_0),
\]
$D(q_0)$ is also nonsquare.  The quadratic field
$\mathbb Q(\sqrt{D(q_0)})$ has a nontrivial automorphism
which exchanges the two numbers in (P-F66b).  Their
coefficients of $\sqrt{D(q_0)}$ are nonzero, since
$F(q_0)\ne0$.  Hence they both have degree two.

For $b\ne0$ rational and $r\ne0$ integral, this same
automorphism exchanges the two numbers in (P-F66c).
Their complex absolute values are
\[
                 |b|\alpha_+(q_0)^r
                 \quad\hbox{and}\quad
                 |b|\alpha_-(q_0)^r,
                                                        \tag{P-F66g}
\]
which are different because the two positive bases are
different and the map $x\mapsto x^r$ is injective on
$\mathbb R_{>0}$.  In particular the numbers remain
quadratic; they cannot become rational after taking
these powers.  Their conjugates do not have a common
absolute value, so no weight $w$ and no prime power $Q$
can satisfy the Weil condition.

For the cohomological consequence we use the known
weight theorem, not a new purity assertion.  Deligne,
\emph{La conjecture de Weil. II}, Definitions~1.2.1--1.2.2,
printed p.~153, defines purity using all complex
conjugates and mixedness by a finite filtration with
pure successive quotients.  Variants~1.2.3--1.2.4,
p.~154, apply this to Frobenius modules.  Theorem~3.3.1,
p.~204, says that compactly supported direct images
preserve mixedness, with the indicated weight bounds.
Corollary~3.3.4, p.~206, is the direct cohomological
statement used here: each Frobenius eigenvalue is algebraic
and all its complex conjugates have absolute value
$Q^{w/2}$ for an integer $w$:
\url{https://www.numdam.org/item/PMIHES_1980__52__137_0.pdf}.
Its introduction, p.~138, explicitly spells out the
specialization to a point: each cohomological
Frobenius eigenvalue has an integer weight and all
its complex conjugates have the corresponding common
absolute value.

Choose the largest of the finitely many weights in a mixed
filtration of $\mathcal L$, and apply that corollary.
Each eigenvalue in (P-F66d) is consequently a Weil number
of some weight.  A stable subquotient has eigenvalues
among those of the original Frobenius module, as is
seen by writing its matrix in block upper-triangular
form.  Finite direct sums take unions of spectra.
A Tate twist multiplies every eigenvalue by an integral
power of the rational number $Q$ and preserves the
all-conjugates condition.  The unequal absolute values
in (P-F66g) exclude every listed possibility.
\end{proof}

\paragraph{Normalizations, finite extensions, and the precise remaining dictionary.}
The assertion already covers any common rational-function
normalization $J(q)\in\mathbb Q(q)$ at a point $q_0$ where
it is defined and nonzero: set $b=J(q_0)$.
It covers source sign changes, Tate twists and finite
base-field extension, since these multiply by integral
powers of the field size or replace Frobenius eigenvalues
by positive integral powers.  Every finite power regrading
$q=t^k$, $k\ge1$, is covered on the sequence $t=3^e$, because
$q_0=3^{ek}$ still satisfies (P-F66f).

The all-complex-conjugates hypothesis is essential to the
stated use of the weight theorem.  We do not silently replace
it by an arbitrary fixed-$\iota$ notion with unconstrained
real weights, and we make no assertion about arbitrary
Weil sheaves lacking the geometric mixedness hypothesis.

This obstruction is different from (P-F52w): it already
excludes one proposed Frobenius eigenvalue at each of the
specified fibers, without requiring both determinant
coefficients to be traces of the same operator over every
finite extension.  It does not exclude a trace of a larger
cohomological representation, or an eigenvalue of an
additional correspondence or auxiliary endomorphism acting
on cohomology.  Such an auxiliary operator need not itself
be Frobenius and is not subject to the Weil-number condition.
In particular the actual positive central-word operator,
the separate Tate--Lefschetz realizations of $T_1,N_1$ in
P-F39, and the completed Euler realization in P-F52 remain
valid.  A comparison using trace or local-system data must
still specify which operator is Frobenius and how the
actual nonstandard source word and its prescribed
normalization enter that operator.
